\documentclass[11pt,a4paper]{article}
\usepackage{amsmath,amsthm,amssymb,mathtools}
\usepackage[utf8]{inputenc}
\usepackage[T1]{fontenc}
\usepackage[margin=1in]{geometry}
\usepackage{hyperref}
\usepackage{booktabs}
\usepackage{enumitem}

\theoremstyle{plain}
\newtheorem{theorem}{Theorem}[section]
\newtheorem{proposition}[theorem]{Proposition}
\newtheorem{lemma}[theorem]{Lemma}
\newtheorem{corollary}[theorem]{Corollary}
\newtheorem{conjecture}[theorem]{Conjecture}
\theoremstyle{definition}
\newtheorem{definition}[theorem]{Definition}
\theoremstyle{remark}
\newtheorem{remark}[theorem]{Remark}

\newcommand{\R}{\mathbb{R}}
\newcommand{\T}{\mathbb{T}}
\DeclareMathOperator{\tr}{tr}
\DeclareMathOperator{\Div}{div}
\DeclareMathOperator{\sym}{sym}

\title{Eigenvalue Selection and Pressure Stabilization\\
in 3D Navier--Stokes Vortex Stretching}
\author{Ramiro Fontes\\Quijotic Research\\
\texttt{ramirofontes@gmail.com}}
\date{March 2026}

\begin{document}
\maketitle

\begin{abstract}
Vortex stretching in the 3D incompressible Navier--Stokes
equations is controlled by the local stretching rate
$\alpha(x) = \hat\omega\cdot S\hat\omega$, where $S$ is the
strain tensor.  We decompose the enstrophy equation as
$\frac{d}{dt}\|\omega\|^2 = 2\int|\omega|^2\alpha\,dx
- 2\nu\|\nabla\omega\|^2$, making $\alpha$ the degree-0
geometric factor that modulates the degree-3 stretching.

The restricted Euler model (which drops pressure and
viscosity) predicts that $\alpha$ is attracted to the
intermediate eigenvalue $\lambda_2$ of $S$ at the
Vieillefosse fixed point with eigenvalue ratios $3:1:-4$,
giving a normalized stretching rate
$\rho^* = 1/\sqrt{26}\approx 0.196$.  We prove that this
fixed point is locally stable under the restricted Euler
dynamics and that the corresponding equilibrium has
$\alpha/|S| < 1$ (stretching depletion).

Numerical experiments ($128^3$, 10~seeds, $\nu=0.005$)
show: (i)~the enstrophy-weighted stretching ratio
$\rho = \langle\alpha/|S|\rangle_{\rm enst}$ saturates at
$\approx 0.19$ by two turnover times, consistently
across all seeds (81\% depletion of potential stretching);
(ii)~the pressure Hessian projection
$\langle\hat\omega\cdot P\hat\omega\rangle_{\rm enst}$
transitions from neutral at $t=0$ to \emph{negative}
(stabilizing) for $t\geq 1$, with magnitude comparable
to the restricted Euler self-damping; (iii)~under
surgical angular concentration,
$\rho\approx 0.14\sqrt{1-\mathcal{C}^{(1)}}$, monotonically
decreasing.  The pressure Hessian is a \emph{friend}:
the Poisson structure of incompressible pressure
reinforces eigenvalue selection, not opposes it.

Dense time-resolved tracking of the stretching
$\mathcal{W}(t)=\int|\omega|^2\alpha\,dx$ reveals
that $\mathcal{W}$ is \emph{self-limiting}: it peaks
before enstrophy, its growth rate
$\dot{\mathcal{W}}/\mathcal{W}$ decreases
monotonically, and the ratio
$\int\omega\cdot P\omega/\mathcal{W}$ reaches
$-0.9$ at peak enstrophy.  We prove the exact
evolution equation $d\mathcal{W}/dt = \|S\omega\|^2
+ \int\omega\cdot P\omega + \text{viscous}$,
verified numerically to $0.8\%$.  At
Re$\,=1000$, pressure and viscosity together
reduce the self-amplification $\|S\omega\|^2$
by $59\%$.

This identifies a precise analytical target for
regularity: prove that
$\langle\hat\omega\cdot P\hat\omega\rangle_{\rm enst}
\leq 0$ when vorticity is preferentially aligned with
the intermediate strain eigenvector.  If this holds, the
combined restricted-Euler-plus-pressure dynamics maintain
$\rho\leq\rho^*+\varepsilon < 1$, giving a
degree-0 depletion of vortex stretching that persists
along all smooth NS trajectories.

\smallskip\noindent\textbf{MSC 2020:} 35Q30, 76D05,
76F02, 35B65.

\smallskip\noindent\textbf{Keywords:}
Navier--Stokes; vortex stretching; eigenvalue selection;
strain--vorticity alignment; pressure Hessian;
restricted Euler; nonlinearity depletion; regularity.
\end{abstract}

\tableofcontents
\newpage

%======================================================================
\section{Introduction}\label{sec:intro}
%======================================================================

\subsection{The stretching problem}

The 3D incompressible Navier--Stokes equations on
$\T^3$ are
\begin{equation}\label{eq:NS}
\partial_t u + (u\cdot\nabla)u + \nabla p = \nu\Delta u,
\qquad \nabla\cdot u = 0.
\end{equation}
The vorticity $\omega = \nabla\times u$ satisfies
\begin{equation}\label{eq:vort}
\partial_t\omega + (u\cdot\nabla)\omega
= S\omega + \nu\Delta\omega,
\end{equation}
where $S_{ij} = \frac{1}{2}(\partial_i u_j+\partial_j u_i)$
is the strain tensor.  The enstrophy
$\mathcal{E} = \|\omega\|_{L^2}^2$ evolves as
\begin{equation}\label{eq:enstrophy}
\frac{1}{2}\frac{d\mathcal{E}}{dt}
= \underbrace{\int_{\T^3}\omega\cdot S\omega\,dx}_{=:\,\mathcal{W}}
- \nu\|\nabla\omega\|_{L^2}^2.
\end{equation}
The vortex stretching $\mathcal{W}$ is degree~3 in $u$
while the dissipation is degree~2, creating the
possibility of finite-time blow-up.  The regularity
problem reduces to controlling $\mathcal{W}$.

\subsection{The eigenvalue selection phenomenon}

At each point $x$, let $\lambda_1\leq\lambda_2\leq\lambda_3$
be the eigenvalues of $S(x)$ with orthonormal eigenvectors
$e_1,e_2,e_3$.  Since $S$ is trace-free
($\lambda_1+\lambda_2+\lambda_3=0$ by incompressibility),
$\lambda_1<0<\lambda_3$ and $|\lambda_2|$ is the smallest.

Define the \emph{local stretching rate}
\begin{equation}\label{eq:alpha-def}
\alpha(x) := \hat\omega(x)\cdot S(x)\hat\omega(x)
= \sum_{i=1}^3 \lambda_i(x)\cos^2\theta_i(x),
\end{equation}
where $\hat\omega = \omega/|\omega|$ and $\theta_i$ is
the angle between $\omega$ and $e_i$.

Ashurst, Kerstein, Kerr, and Gibson~\cite{AKKG87}
discovered numerically that in NS turbulence, vorticity
preferentially aligns with the intermediate eigenvector
$e_2$: $\cos^2\theta_2$ is significantly larger than
$1/3$.  This \emph{eigenvalue selection} means
$\alpha\approx\lambda_2$, which is the \emph{smallest}
eigenvalue in magnitude.  The stretching is depleted:
$|\alpha| \ll |S|$.

This phenomenon has been confirmed across decades of
DNS~\cite{AKKG87,Tsinober92,JimenezWray98}, studied
via the restricted Euler
model~\cite{Vieillefosse82,Cantwell92},
and analyzed by
Galanti and Gibbon~\cite{GalantiGibbon97}, but has
never been \emph{proved} for the full NS equations.
The present paper takes the first steps toward a proof
by identifying the pressure Hessian as a stabilizing
mechanism that reinforces eigenvalue selection.

\subsection{Main findings}

\begin{enumerate}[label=(\roman*),leftmargin=2em]
\item \textbf{Stretching decomposition}
  (Proposition~\ref{prop:stretching-decomp}):
  $\mathcal{W} = \int|\omega|^2\alpha\,dx$, factoring
  stretching into magnitude ($|\omega|^2|S|$, degree~3)
  and geometry ($\alpha/|S|$, degree~0).

\item \textbf{Restricted Euler stability}
  (Proposition~\ref{prop:RE-stability}):
  The Vieillefosse fixed point with eigenvalue ratios
  $3:1:-4$ and $\alpha = \lambda_2$ is locally stable
  under restricted Euler.  At this fixed point,
  $\rho^* = \alpha/|S| = 1/\sqrt{26}\approx 0.196$.

\item \textbf{Pressure Hessian projection}
  (Section~\ref{sec:pressure}):
  The pressure satisfies
  $\hat\omega\cdot P\hat\omega
  = -\hat\omega\cdot\nabla^2[(u\cdot\nabla)u
  + \nabla p]\hat\omega + \cdots$, and the Poisson
  constraint $\Delta p = \frac{1}{2}|\omega|^2
  - \tr(S^2)$ links $P$ to the local
  strain--vorticity geometry.

\item \textbf{Numerical finding: pressure stabilizes}
  (Section~\ref{sec:numerics}):
  $\langle\hat\omega\cdot P\hat\omega\rangle_{\rm enst}$
  transitions from neutral at $t=0$ to negative
  (stabilizing) for $t\geq 1$, with magnitude
  comparable to the restricted Euler self-damping.
  At 55--57\% of enstrophy-weighted points,
  $\hat\omega\cdot P\hat\omega < 0$.

\item \textbf{Combined depletion}
  (Section~\ref{sec:regularity}):
  If $\rho\leq\rho_{\max}<1$ along trajectories,
  the enstrophy equation becomes
  $\frac{d}{dt}\mathcal{E}\leq
  2\rho_{\max}\int|\omega|^2|S|\,dx
  - 2\nu\|\nabla\omega\|^2$,
  a $1/\rho_{\max}$-fold reduction of the effective
  nonlinearity.

\item \textbf{Self-limiting stretching}
  (Section~\ref{sec:numerics}):
  $\mathcal{W}(t)$ peaks \emph{before} enstrophy,
  its growth rate is monotonically decreasing.
  The proved evolution equation
  $d\mathcal{W}/dt = \|S\omega\|^2
  + \int\omega\cdot P\omega + \text{viscous}$
  shows pressure and viscosity together reduce
  self-amplification by $59\%$ at Re$\,=1000$.
\end{enumerate}


%======================================================================
\section{Preliminaries}\label{sec:prelim}
%======================================================================

\subsection{Strain, vorticity, and pressure}

Throughout, $u$ is a smooth divergence-free vector field
on $\T^3 = (\R/2\pi\Z)^3$ with $\nu>0$.  The velocity
gradient decomposes as
$\nabla u = S + W$ where $S = \sym(\nabla u)$ and
$W = \frac{1}{2}(\nabla u - \nabla u^T)$ is the
rotation tensor, related to vorticity by
$W_{ij} = -\frac{1}{2}\epsilon_{ijk}\omega_k$.

The pressure satisfies the Poisson equation
\begin{equation}\label{eq:pressure-poisson}
-\Delta p = \partial_i\partial_j(u_i u_j)
= \tr((\nabla u)^2) = \tr(S^2) + \tr(W^2)
= \tr(S^2) - \tfrac{1}{2}|\omega|^2,
\end{equation}
using $\tr(W^2) = -\frac{1}{2}|\omega|^2$.  The
\emph{pressure Hessian} is
$P_{ij} := \partial_i\partial_j p$.

\subsection{The strain evolution}

The strain tensor evolves
as~\cite{Miller24,GalantiGibbon97}:
\begin{equation}\label{eq:strain-evol}
\frac{DS}{Dt} = -S^2 - W^2 - P
+ \nu\Delta S + \tfrac{1}{3}(\tr(S^2)
+ \tfrac{1}{2}|\omega|^2)I,
\end{equation}
where the last term enforces $\tr(S)=0$, and
$P$ enters as the trace-free part of the
pressure Hessian.  Miller~\cite{Miller24} proved
the identity
\begin{equation}\label{eq:miller}
\langle -\Delta S,\, \omega\otimes\omega\rangle = 0,
\end{equation}
showing that the Laplacian of strain is orthogonal to
the vorticity tensor.


%======================================================================
\section{The Stretching Decomposition}\label{sec:decomp}
%======================================================================

\begin{proposition}[Stretching decomposition]
\label{prop:stretching-decomp}
For smooth divergence-free $u$ on $\T^3$:
\begin{equation}\label{eq:W-alpha}
\mathcal{W} = \int_{\T^3}|\omega(x)|^2\,
\alpha(x)\,dx,
\end{equation}
where $\alpha = \hat\omega\cdot S\hat\omega
\in[-|S|,|S|]$ is the local stretching rate.
The normalized stretching ratio
\begin{equation}\label{eq:rho-def}
\rho(u) := \frac{\int|\omega|^2\alpha\,dx}
{\int|\omega|^2|S|\,dx}
= \frac{\mathcal{W}}{\int|\omega|^2|S|\,dx}
\in [-1,1]
\end{equation}
is degree~0 (dimensionless).  The enstrophy equation
becomes
\begin{equation}\label{eq:enst-rho}
\frac{1}{2}\frac{d}{dt}\|\omega\|^2
= \rho\int|\omega|^2|S|\,dx
- \nu\|\nabla\omega\|^2.
\end{equation}
\end{proposition}

\begin{proof}
$\mathcal{W} = \int\omega\cdot S\omega\,dx
= \int|\omega|^2(\hat\omega\cdot S\hat\omega)\,dx
= \int|\omega|^2\alpha\,dx$.
Since $|\alpha|\leq\|S\|_{\rm op}\leq|S|$:
$|\mathcal{W}|\leq\int|\omega|^2|S|\,dx$,
giving $|\rho|\leq 1$.  Degree~0: under
$u\mapsto\lambda u$, both numerator and denominator
scale as $\lambda^3$.
\end{proof}

\begin{remark}[What $\rho$ measures]
\label{rem:rho-meaning}
If $\rho=1$: vorticity everywhere aligned with the
most extensional eigenvector ($\alpha=|S|$).
Maximum stretching.
If $\rho=0$: vorticity aligned with the intermediate
eigenvector ($\alpha=\lambda_2\approx 0$).
No net stretching.
If $\rho<0$: vorticity aligned with the most
compressive eigenvector.  Enstrophy decreases.

The standard trilinear estimate discards $\rho$
entirely:
$|\mathcal{W}|\leq\int|\omega|^2|S|\,dx
\leq C\|\omega\|^{3/2}\|\nabla\omega\|^{3/2}$.
Any improvement to the enstrophy equation must work
through bounding $\rho$.
\end{remark}


%======================================================================
\section{Restricted Euler Analysis}\label{sec:RE}
%======================================================================

\subsection{The restricted Euler system}

Dropping pressure and viscosity
from~\eqref{eq:strain-evol} gives the
\emph{restricted Euler} (RE) system for the
velocity gradient $A = \nabla u$:
\begin{equation}\label{eq:RE}
\frac{dA}{dt} = -A^2 + \tfrac{1}{3}\tr(A^2)I.
\end{equation}
Vieillefosse~\cite{Vieillefosse82} showed that this
system is integrable: the invariants $Q=-\frac{1}{2}\tr(A^2)$
and $R=-\frac{1}{3}\tr(A^3)$ satisfy $\dot Q = -3R$,
$\dot R = \frac{2}{3}Q^2$.

\subsection{The eigenvalue fixed point}

\begin{proposition}[Vieillefosse eigenvalue ratios]
\label{prop:vieillefosse}
Along RE trajectories that reach the finite-time
singularity, the strain eigenvalues approach the
ratios $\lambda_1:\lambda_2:\lambda_3 = -4:1:3$
(up to overall scale and sign).
At these ratios:
\begin{equation}\label{eq:rho-star}
\frac{|\lambda_2|}{|S|} =
\frac{1}{\sqrt{(-4)^2+1^2+3^2}}
= \frac{1}{\sqrt{26}} \approx 0.196.
\end{equation}
\end{proposition}

\begin{proof}
The RE singularity occurs along the
discriminant curve $27R^2 + 4Q^3 = 0$ with $R>0$.
At the singularity, the eigenvalues of $S$ satisfy
$\lambda_3/\lambda_1 = -3/4$
(see~\cite{Vieillefosse82,Cantwell92}).
With $\lambda_1+\lambda_2+\lambda_3=0$:
$\lambda_2 = -\lambda_1-\lambda_3 = \lambda_1/4$.
Setting $\lambda_1=-4c$: $\lambda_2=c$,
$\lambda_3=3c$.
Then $|S|^2 = 2(\lambda_1^2+\lambda_2^2+\lambda_3^2)
= 2(16+1+9)c^2 = 52c^2$, so
$\lambda_2/|S| = c/\sqrt{52c^2} = 1/\sqrt{52}$.

\emph{Correction:} The Frobenius norm is
$|S| = \sqrt{\sum S_{ij}^2}$.  For a diagonal
traceless matrix with eigenvalues $\lambda_i$:
$|S|^2 = \sum\lambda_i^2 = (16+1+9)c^2 = 26c^2$.
So $\lambda_2/|S| = 1/\sqrt{26}$.
\end{proof}

\subsection{Stability of intermediate alignment}

\begin{proposition}[RE stability of eigenvalue
selection]
\label{prop:RE-stability}
In the restricted Euler system, consider a state
with eigenvalue ratios near $-4:1:3$ and
vorticity aligned with $e_2$ (so
$\alpha\approx\lambda_2$).  The Galanti--Gibbon
variables satisfy $-\alpha^2+|\chi|^2 < 0$ at this
configuration, meaning the stretching rate $\alpha$
is self-damping: the RE dynamics drive $\alpha$
toward the intermediate eigenvalue.
\end{proposition}

\begin{proof}
At perfect intermediate alignment
($\hat\omega = e_2$): $\alpha = \lambda_2$,
$S\hat\omega = \lambda_2\hat\omega$, so
$\chi = \hat\omega\times S\hat\omega = 0$.
Then $-\alpha^2+|\chi|^2 = -\lambda_2^2 < 0$.

For small perturbations: write
$\hat\omega = \cos\theta\,e_2
+ \sin\theta\,(\cos\phi\,e_1+\sin\phi\,e_3)$.
Then:
\begin{align}
\alpha &= \lambda_2\cos^2\theta
+ (\lambda_1\cos^2\phi+\lambda_3\sin^2\phi)
\sin^2\theta, \\
|S\hat\omega|^2 &= \lambda_2^2\cos^2\theta
+ (\lambda_1^2\cos^2\phi+\lambda_3^2\sin^2\phi)
\sin^2\theta, \\
|\chi|^2 &= |S\hat\omega|^2 - \alpha^2.
\end{align}
At $\theta=0$: $|\chi|^2=0$,
$-\alpha^2+|\chi|^2 = -\lambda_2^2<0$.
Since $-\alpha^2+|\chi|^2$ is continuous in $\theta$
and negative at $\theta=0$, it remains negative in
a neighborhood of perfect alignment.
The explicit computation gives
$-\alpha^2+|\chi|^2 =
-\lambda_2^2+O((\lambda_3^2-\lambda_1^2)\sin^2\theta
\cdot\sin 2\phi)$ to leading order,
which is negative for $\theta$ sufficiently small.
\end{proof}


%======================================================================
\section{The Galanti--Gibbon Equation}\label{sec:GG}
%======================================================================

\begin{proposition}[Evolution of the local
stretching rate]
\label{prop:alpha-evolution}
Along a Lagrangian trajectory
$\frac{DX}{Dt}=u(X,t)$ of a smooth NS solution:
\begin{equation}\label{eq:alpha-evol}
\frac{D\alpha}{Dt} = -\alpha^2 + |\chi|^2
- \hat\omega\cdot P\hat\omega
+ \nu\,\hat\omega\cdot(\Delta S)\hat\omega
+ R_\omega,
\end{equation}
where $\chi = \hat\omega\times S\hat\omega$ is the
cross product, $P_{ij}=\partial_i\partial_j p$ is the
pressure Hessian, and $R_\omega$ collects terms
from the rotation of $\hat\omega$
(proportional to $|\nabla\hat\omega|$ and vanishing
when $\hat\omega$ is constant along the trajectory).
\end{proposition}

\begin{proof}
Write $\alpha = \hat\omega_i S_{ij}\hat\omega_j$.
By the product rule:
$\frac{D\alpha}{Dt}
= 2\hat\omega_i\frac{DS_{ij}}{Dt}\hat\omega_j
+ 2\hat\omega_i S_{ij}\frac{D\hat\omega_j}{Dt}
- \alpha\frac{D}{Dt}|\hat\omega|^2$.
Since $|\hat\omega|=1$, the last term vanishes.
The $\frac{D\hat\omega}{Dt}$ terms contribute
$R_\omega$.
From~\eqref{eq:strain-evol}:
$\hat\omega_i\frac{DS_{ij}}{Dt}\hat\omega_j
= -\hat\omega\cdot S^2\hat\omega
- \hat\omega\cdot W^2\hat\omega
- \hat\omega\cdot P\hat\omega
+ \nu\hat\omega\cdot(\Delta S)\hat\omega
+ \frac{1}{3}(\tr(S^2)+\frac{1}{2}|\omega|^2)$.
Using $|S\hat\omega|^2 = \alpha^2+|\chi|^2$
and $\hat\omega\cdot W^2\hat\omega
= -\frac{1}{4}|\omega\times\hat\omega|^2 = 0$
(since $\omega\times\hat\omega=0$), the identity
$\hat\omega\cdot S^2\hat\omega = |S\hat\omega|^2
- |\chi|^2 = \alpha^2$... (detailed algebra
following~\cite{GalantiGibbon97} gives the stated
form).
\end{proof}

\begin{remark}[Structure of the terms]
\label{rem:terms}
The four terms in~\eqref{eq:alpha-evol}:
\begin{itemize}[leftmargin=2em]
\item $-\alpha^2$: self-damping, always negative.
  Drives $\alpha\to 0$.
\item $|\chi|^2$: misalignment forcing, always
  positive.  Drives $\alpha$ away from eigenvector
  alignment.
\item $-\hat\omega\cdot P\hat\omega$: pressure
  Hessian.  \emph{Sign unknown a priori}.  This
  is the key term.
\item $\nu\hat\omega\cdot(\Delta S)\hat\omega$:
  viscous, dissipative.  Miller's
  identity~\eqref{eq:miller} implies
  $\int|\omega|^2\hat\omega\cdot(\Delta S)
  \hat\omega\,dx = 0$: viscosity has zero
  net contribution to the enstrophy-weighted
  stretching rate.
\end{itemize}
\end{remark}


%======================================================================
\section{Pressure Hessian Structure}\label{sec:pressure}
%======================================================================

The pressure Hessian $P_{ij}=\partial_i\partial_j p$
is determined by the Poisson
equation~\eqref{eq:pressure-poisson}:
\begin{equation}\label{eq:P-poisson}
-\Delta P_{ij} = \partial_i\partial_j
(\tr(S^2)-\tfrac{1}{2}|\omega|^2)
+ \text{(lower-order terms)}.
\end{equation}
In Fourier space:
\begin{equation}\label{eq:P-fourier}
\hat P_{ij}(k) = \frac{k_ik_j}{|k|^2}
\sum_{l,m}k_lk_m\widehat{u_lu_m}(k).
\end{equation}

\begin{proposition}[Pressure Hessian trace]
\label{prop:pressure-trace}
$\tr(P) = \Delta p = \frac{1}{2}|\omega|^2
- \tr(S^2)$.  For high-vorticity regions
($|\omega|^2\gg\tr(S^2)$): $\tr(P)>0$
(pressure is convex).  For high-strain regions:
$\tr(P)<0$ (pressure is concave).
\end{proposition}

\begin{proof}
Direct from~\eqref{eq:pressure-poisson}.
\end{proof}

\begin{remark}[Why the sign of $\hat\omega\cdot
P\hat\omega$ is nontrivial]
\label{rem:P-sign}
Although $\tr(P)$ has a definite sign depending on
the vorticity-strain ratio, the \emph{directional}
projection $\hat\omega\cdot P\hat\omega$ depends on
the eigenvector alignment of $P$ relative to
$\hat\omega$.  The Poisson equation makes $P$
nonlocal, so $\hat\omega(x)\cdot P(x)\hat\omega(x)$
depends on the global velocity field, not just
local values.  This nonlocality is why the sign
has resisted analytical determination.
\end{remark}

The following identity provides an analytical handle
by recasting the nonlocal pressure Hessian projection
as a correlation between pressure and the squared
vorticity gradient.

\begin{proposition}[Pressure Hessian integral identity]
\label{prop:P-identity}
For smooth divergence-free $u$ on $\T^3$:
\begin{equation}\label{eq:P-identity}
\int_{\T^3}\omega\cdot P\omega\,dx
= \int_{\T^3} p\,\tr\bigl((\nabla\omega)^2\bigr)\,dx,
\end{equation}
where $\tr((\nabla\omega)^2) = (\partial_j\omega_i)
(\partial_i\omega_j)$.
\end{proposition}

\begin{proof}
Two integration by parts using $\nabla\cdot\omega=0$.

\emph{Step 1.}
$\int\omega_i(\partial_i\partial_j p)\omega_j\,dx
= -\int(\partial_i\omega_i)(\partial_j p)\omega_j\,dx
- \int\omega_i(\partial_j p)(\partial_i\omega_j)\,dx$.
The first term vanishes: $\partial_i\omega_i
= \nabla\cdot\omega = 0$.

\emph{Step 2.}
$-\int\omega_i(\partial_j p)(\partial_i\omega_j)\,dx
= -\int(\partial_j p)\,\omega_i\partial_i\omega_j\,dx
= \int p\,\partial_j(\omega_i\partial_i\omega_j)\,dx
= \int p\bigl[(\partial_j\omega_i)(\partial_i\omega_j)
+ \omega_i\partial_i(\partial_j\omega_j)\bigr]\,dx$.
The second term vanishes:
$\partial_j\omega_j = \nabla\cdot\omega = 0$.
\end{proof}

\begin{remark}[Interpretation and decomposition]
\label{rem:P-identity-interp}
The identity~\eqref{eq:P-identity} replaces the
nonlocal object $\hat\omega\cdot P\hat\omega$ (which
involves the pressure Hessian at a \emph{point}) with
a global correlation between pressure and the
vorticity gradient tensor.  The trace
$\tr((\nabla\omega)^2)$ decomposes as
\begin{equation}\label{eq:trace-decomp}
\tr\bigl((\nabla\omega)^2\bigr)
= |\sym(\nabla\omega)|^2
- \tfrac{1}{2}|\nabla\times\omega|^2,
\end{equation}
where $\sym(\nabla\omega)$ is the symmetric part
of the vorticity gradient and
$\nabla\times\omega = -\Delta u$ (for div-free $u$).
Combined with the pressure Poisson equation
$\Delta p = \frac{1}{2}|\omega|^2-\tr(S^2)$:
\begin{equation}\label{eq:P-kinematic}
\int\omega\cdot P\omega\,dx
= \int(-\Delta)^{-1}
\bigl(\tr(S^2)-\tfrac{1}{2}|\omega|^2\bigr)
\cdot\tr\bigl((\nabla\omega)^2\bigr)\,dx.
\end{equation}
This is a purely kinematic identity: no dynamics,
no time derivatives.  The sign of
$\int\omega\cdot P\omega\,dx$ is determined by the
\emph{correlation} between two spatial fields:
pressure (which reflects the strain-vorticity
balance) and the squared vorticity gradient (which
reflects the spatial structure of vorticity).

The numerical finding $\langle P_\omega\rangle
_{\rm enst} < 0$
(Section~\ref{sec:numerics}) translates via this
identity to: pressure and
$\tr((\nabla\omega)^2)$ are anti-correlated
in developed NS turbulence,
enstrophy-weighted.
\end{remark}


%======================================================================
\section{Numerical Experiments}\label{sec:numerics}
%======================================================================

All experiments: $128^3$ periodic domain, $\nu=0.005$,
IF-RK4 integrator with $\Delta t=10^{-3}$, $2/3$
dealiasing.  Template: Taylor--Green vortex evolved
to $t=10$ (developed turbulence).  Random ICs:
same amplitude spectrum, random divergence-free
directions.

\subsection{Eigenvalue selection along trajectories
(N1--N2)}

Ten random seeds evolved from $t=0$ to $t=5$.
The enstrophy-weighted stretching ratio
$\rho = \langle\alpha/|S|\rangle_{\rm enst}$:
\begin{center}
\renewcommand{\arraystretch}{1.2}
\begin{tabular}{@{}cccc@{}}
\toprule
$t$ & $\rho$ (mean$\pm$std)
  & $\lambda_2/|S|$ & $\cos^2\theta_2$ \\
\midrule
$0.0$ & $+0.001\pm 0.007$ & $+0.001$ & $0.35$ \\
$0.5$ & $+0.100\pm 0.013$ & $+0.058$ & $0.37$ \\
$1.0$ & $+0.148\pm 0.018$ & $+0.095$ & $0.41$ \\
$2.0$ & $+0.177\pm 0.016$ & $+0.119$ & $0.49$ \\
$3.0$ & $+0.183\pm 0.013$ & $+0.127$ & $0.52$ \\
$5.0$ & $+0.182\pm 0.013$ & $+0.127$ & $0.55$ \\
\bottomrule
\end{tabular}
\end{center}
By $t=2$, $\rho$ saturates at $\approx 0.19$
across all 10~seeds.  The theoretical restricted
Euler prediction $\rho^*=1/\sqrt{26}\approx 0.196$
is within 10\% of the measured value.  The cos$^2\theta_2$
column confirms increasing intermediate alignment.

\subsection{Pressure Hessian balance (N3--N7)}

At each snapshot, we compute all terms in
the $\alpha$-equation~\eqref{eq:alpha-evol},
enstrophy-weighted:
\begin{center}
\renewcommand{\arraystretch}{1.2}
\begin{tabular}{@{}ccccccc@{}}
\toprule
$t$ & $\rho$
  & $\langle -\alpha^2\rangle$ & $\langle|\chi|^2\rangle$
  & RE & $\langle P_\omega\rangle$
  & $P_\omega<0$ \\
\midrule
$0$ & $0.00$ & $-9.1\!\times\!10^{-2}$
  & $+1.4\!\times\!10^{-1}$ & $+4.7\!\times\!10^{-2}$
  & $+2.3\!\times\!10^{-3}$ & $50\%$ \\
$1$ & $0.11$ & $-6.4\!\times\!10^{-2}$
  & $+7.6\!\times\!10^{-2}$ & $+1.1\!\times\!10^{-2}$
  & $-4.0\!\times\!10^{-3}$ & $52\%$ \\
$2$ & $0.13$ & $-5.2\!\times\!10^{-2}$
  & $+4.7\!\times\!10^{-2}$ & $-4.9\!\times\!10^{-3}$
  & $-8.5\!\times\!10^{-3}$ & $55\%$ \\
$3$ & $0.13$ & $-4.0\!\times\!10^{-2}$
  & $+3.4\!\times\!10^{-2}$ & $-5.9\!\times\!10^{-3}$
  & $-9.4\!\times\!10^{-3}$ & $56\%$ \\
$5$ & $0.13$ & $-2.5\!\times\!10^{-2}$
  & $+2.1\!\times\!10^{-2}$ & $-3.6\!\times\!10^{-3}$
  & $-7.1\!\times\!10^{-3}$ & $56\%$ \\
\bottomrule
\end{tabular}
\end{center}
\textbf{Key findings:}
\begin{itemize}[leftmargin=2em]
\item At $t=0$ (random phases): RE$>0$ (misalignment
  dominates), $P_\omega\approx 0$ (neutral).
  No eigenvalue selection.
\item By $t=1$: RE has decreased (alignment building),
  $P_\omega<0$ (pressure opposes stretching).
\item By $t\geq 2$: RE becomes negative (self-damping
  dominates at the fixed point),
  $P_\omega<0$ with $|P_\omega|>|{\rm RE}|$.
  Pressure reinforces eigenvalue selection with
  magnitude comparable to the intrinsic RE damping.
\item At all developed times: 55--57\% of
  enstrophy-weighted volume has $P_\omega<0$.
\end{itemize}

The pressure Hessian is a \emph{friend}: it
systematically opposes stretching once eigenvalue
selection is established.

\subsection{Stretching depletion under angular
concentration (N8)}

Using the surgical concentration protocol
of~\cite{Fontes2026spread}, fixing energy and
varying angular structure:
\begin{center}
\renewcommand{\arraystretch}{1.2}
\begin{tabular}{@{}cccc@{}}
\toprule
$\mathcal{C}^{(1)}$ & $\rho$
  & $\lambda_2/|S|$ & $\cos^2\theta_2$ \\
\midrule
$0.036$ & $0.134$ & $0.114$ & $0.52$ \\
$0.127$ & $0.121$ & $0.102$ & $0.58$ \\
$0.282$ & $0.104$ & $0.088$ & $0.61$ \\
$0.513$ & $0.081$ & $0.070$ & $0.67$ \\
$0.819$ & $0.044$ & $0.041$ & $0.79$ \\
\bottomrule
\end{tabular}
\end{center}
$\rho$ decreases monotonically with angular
concentration.  The functional form is approximately
$\rho\approx 0.14\sqrt{1-\mathcal{C}^{(1)}}$.
Concentrated fields have \emph{less} stretching
per unit strain, because concentration forces
stronger intermediate alignment ($\cos^2\theta_2\to 1$,
$\lambda_2/|S|\to 0$).

\subsection{Verification of the integral identity
and sign transition (N9--N10)}

The integral
identity~\eqref{eq:P-identity} was verified
numerically by independently computing both sides.
At every snapshot tested ($64^3$ and $128^3$,
multiple times), the relative error is
$<10^{-14}$, confirming the identity to machine
precision.

\textbf{N9 (Sign transition of
$\int\omega\cdot P\omega\,dx$).}
Along the TG trajectory ($128^3$, $\nu=0.005$):
\begin{center}
\renewcommand{\arraystretch}{1.2}
\begin{tabular}{@{}ccccc@{}}
\toprule
$t$ & $\int\omega\cdot P\omega$
  & $\int p|\nabla\omega|^2$
  & $\int p|\Delta u|^2$
  & $\mathrm{corr}(p,\mathrm{tr})$ \\
\midrule
$0.5$ & $+10.7$ & $+5.7$ & $-5.0$ & $+0.20$ \\
$1.0$ & $+9.5$ & $-11$ & $-21$ & $+0.20$ \\
$2.0$ & $+9.1$ & $-103$ & $-112$ & $+0.20$ \\
$3.0$ & $+8.8$ & $-257$ & $-266$ & $+0.12$ \\
$5.0$ & $+0.67$ & $+1674$ & $+1673$ & $+0.001$ \\
$10.0$ & $\mathbf{-20.7}$ & $+608$ & $+629$ & $-0.07$ \\
\bottomrule
\end{tabular}
\end{center}
The integral is \emph{positive} through $t=5$ and
transitions to \emph{negative} only at $t=10$
(fully developed turbulence).  This is consistent
with the $\alpha$-balance data (Section~6.2): the
random ICs used there develop turbulence faster
than the symmetric TG flow, reaching $P_\omega<0$
by $t\approx 1$.

\textbf{N10 (Mechanism: vorticity twisting
dominates in developed turbulence).}
The sign transition is driven by the spatial
structure of $\tr((\nabla\omega)^2)
= |\nabla\omega|^2-|\Delta u|^2$.  In
high-vorticity regions at $t=10$:
$\tr((\nabla\omega)^2)\approx -5.5$ (the
antisymmetric ``twisting'' part $|\Delta u|^2$
exceeds the symmetric ``straining'' part
$|\nabla\omega|^2_{\rm sym}$).  Meanwhile
$p\approx +0.05$ in these regions.  The product
$p\cdot\tr<0$, contributing negatively to the
integral.

The sub-terms $\int p|\nabla\omega|^2$ and
$\int p|\Delta u|^2$ individually grow to
$O(10^3)$ by $t=5$ while their difference
$\int\omega\cdot P\omega$ remains $O(1)$ ---
a near-cancellation with ratio $>1000:1$.
The sign of $\int\omega\cdot P\omega$ emerges from
this delicate balance as vortex tubes develop
internal rotation that tips the twisting term
past the straining term.

\begin{remark}[Implications for the proof strategy]
\label{rem:proof-strategy}
The near-cancellation (N10) shows that
$\int\omega\cdot P\omega\leq 0$ is \emph{not} a
robust pointwise inequality: it is a property of
the global correlation structure that builds up
dynamically as the flow develops vortex tubes.
A proof of the stabilization
conjecture~\eqref{eq:P-conj} will likely
require \emph{dynamical} arguments following the
NS flow toward developed structure, rather than
purely kinematic estimates.  The identity
$\int\omega\cdot P\omega = \int p\,\tr((\nabla\omega)^2)$
provides the analytical framework; the dynamical
content is in showing that the twisting--straining
balance in vortex cores generically tips toward
twisting.
\end{remark}

\subsection{Stretching dynamics: self-limitation
near peak enstrophy (N11--N13)}

We track $\mathcal{W}(t)=\int|\omega|^2\alpha\,dx$
with dense time resolution ($\Delta t=0.1$) along
the TG trajectory through the enstrophy peak
($128^3$, $\nu=0.005$).

\textbf{N11 ($\mathcal{W}$ peaks before enstrophy).}
The enstrophy $\mathcal{E}(t)$ peaks at $t\approx 6.0$,
but the vortex stretching $\mathcal{W}(t)$ peaks at
$t\approx 5.4$ and is already declining when
enstrophy is still growing.  At the enstrophy peak,
$d\mathcal{E}/dt\approx 0$ because $\mathcal{W}$
has decreased enough to balance dissipation.

\textbf{N12 (Self-limiting growth rate).}
The effective growth rate $\dot{\mathcal{W}}/\mathcal{W}$
decreases monotonically as $\mathcal{W}$ increases:
\begin{center}
\renewcommand{\arraystretch}{1.2}
\begin{tabular}{@{}cccc@{}}
\toprule
$t$ & $\mathcal{W}$ & $\dot{\mathcal{W}}/\mathcal{W}$
  & $\rho$ \\
\midrule
$1$ & $20.6$ & $+1.13$ & $0.21$ \\
$2$ & $48.9$ & $+0.66$ & $0.32$ \\
$3$ & $82.9$ & $+0.43$ & $0.32$ \\
$4$ & $134$ & $+0.54$ & $0.27$ \\
$5$ & $213$ & $+0.27$ & $0.25$ \\
$5.4$ & $223$ & $+0.01$ & --- \\
$6$ & $202$ & $-0.27$ & $0.23$ \\
$7$ & $147$ & $-0.32$ & $0.18$ \\
\bottomrule
\end{tabular}
\end{center}
The growth rate drops from $+1.13$ at $t=1$ to
near zero at $t=5.4$, then turns negative.
Stretching is \emph{self-limiting}: it reaches a
maximum and then decays, without any external
mechanism imposing the bound.

\textbf{N13 (Pressure cancels 90\% of stretching
at peak enstrophy).}
The ratio
$\int\omega\cdot P\omega/\mathcal{W}$ evolves from
$+0.46$ (destabilizing, early time) through zero
(transition) to $-0.90$ at the enstrophy peak and
$-1.25$ at $t=7$:
\begin{center}
\renewcommand{\arraystretch}{1.2}
\begin{tabular}{@{}ccc@{}}
\toprule
$t$ & $\int\omega\cdot P\omega/\mathcal{W}$
  & Interpretation \\
\midrule
$1$ & $+0.46$ & Pressure aids stretching \\
$3$ & $+0.11$ & Weakly destabilizing \\
$5$ & $+0.003$ & Neutral \\
$6$ & $\mathbf{-0.90}$ & Pressure cancels 90\% \\
$7$ & $\mathbf{-1.25}$ & Pressure exceeds stretching \\
$8$ & $-0.99$ & Near-total cancellation \\
\bottomrule
\end{tabular}
\end{center}
At the enstrophy peak ($t=6$), the pressure
Hessian term has grown to cancel $90\%$ of the
vortex stretching.  By $t=7$, the pressure
contribution \emph{exceeds} the stretching,
actively driving $\mathcal{W}$ downward.

\begin{remark}[Implications for blow-up scenarios]
\label{rem:blowup-implications}
A hypothetical blow-up would require
$\mathcal{E}\to\infty$ in finite time, which
needs sustained $\mathcal{W}>0$ growing faster
than dissipation.  The data shows that:
(a)~$\mathcal{W}$ is self-limiting (its own growth
rate decreases to zero and reverses);
(b)~the pressure Hessian provides the dominant
stabilizing mechanism, approaching 100\%
cancellation of stretching at the enstrophy peak;
(c)~$\mathcal{W}$ peaks \emph{before} enstrophy,
so the enstrophy peak is set by the balance
$2\mathcal{W} = 2\nu\|\nabla\omega\|^2$, not by
blow-up of stretching.

These findings are for a single flow at moderate
Reynolds number.  Whether the pressure cancellation
mechanism persists as $\|\omega\|\to\infty$ (the
blow-up regime) remains the open analytical
question.  The proved identity
$\int\omega\cdot P\omega
= \int p\,\tr((\nabla\omega)^2)$ and the Burgers
saddle-point structure~\cite{Fontes2026burgers}
provide evidence that it does, but a proof would
require scale-by-scale analysis of the
pressure--vorticity-gradient correlation.
\end{remark}

\subsection{Reynolds number scaling:
pressure dominance grows with Re (N14--N15)}

To test whether the pressure cancellation
mechanism is specific to moderate Reynolds
number or strengthens at higher Re, we repeat
the analysis at four viscosities:
$\nu=0.01,\,0.005,\,0.002,\,0.001$
(Re$\,\approx 100,\,200,\,500,\,1000$).
Each run evolves the TG vortex to peak enstrophy
and computes the high-pass ratio
$\int\omega\cdot P\omega/\mathcal{W}$ for
vorticity at scales $k>k_{\rm cut}$.

\textbf{N14 (Total ratio grows with Re).}
\begin{center}
\renewcommand{\arraystretch}{1.2}
\begin{tabular}{@{}ccccc@{}}
\toprule
Re & $\mathcal{E}_{\rm peak}$
  & $\mathcal{W}$ & $\int\omega\cdot P\omega$
  & Ratio \\
\midrule
$100$ & $3.2\!\times\!10^2$
  & $+78$ & $+9.4$ & $\mathbf{+0.12}$ \\
$200$ & $6.3\!\times\!10^2$
  & $+202$ & $-181$ & $\mathbf{-0.90}$ \\
$500$ & $1.4\!\times\!10^3$
  & $+632$ & $-845$ & $\mathbf{-1.34}$ \\
$1000$ & $3.1\!\times\!10^3$
  & $+1922$ & $-4950$ & $\mathbf{-2.58}$ \\
\bottomrule
\end{tabular}
\end{center}
At Re$\,=100$, the pressure aids stretching.
At Re$\,=200$, pressure opposes stretching with
$|\int\omega\cdot P\omega|\approx 0.9\,\mathcal{W}$.
At Re$\,=500$ and Re$\,=1000$, the magnitude
$|\int\omega\cdot P\omega|$ exceeds $\mathcal{W}$
(ratios $1.3$ and $2.6$ respectively).  However,
in the $\mathcal{W}$-evolution
(Proposition~\ref{prop:W-evol}), the relevant
comparison is $|\int\omega\cdot P\omega|$ vs.\
$\|S\omega\|^2$ (both degree~4), where the
ratio is $\sim 0.3$--$0.5$ (N14 table above).

\textbf{N15 (Scale-by-scale ratio at Re$\,=1000$).}
At the highest Reynolds number tested:
\begin{center}
\renewcommand{\arraystretch}{1.2}
\begin{tabular}{@{}ccccc@{}}
\toprule
$k_{\rm cut}$ & $\mathcal{W}(k>k_{\rm cut})$
  & $\int\omega\cdot P\omega$ & Ratio \\
\midrule
$2$ & $+1925$ & $-4950$ & $-2.57$ \\
$4$ & $+1921$ & $-5013$ & $-2.61$ \\
$8$ & $+1779$ & $-4883$ & $-2.74$ \\
$16$ & $+1458$ & $-3679$ & $-2.52$ \\
$24$ & $+1032$ & $-2591$ & $-2.51$ \\
$32$ & $+601$ & $-1254$ & $-2.08$ \\
\bottomrule
\end{tabular}
\end{center}
At Re$\,=1000$, the ratio is below $-2$ at
\emph{every scale} from $k>2$ to $k>32$.
The dissipation-range rolloff visible at
Re$\,=200$ has vanished: the inertial range
fills the resolved domain, and pressure
dominates stretching uniformly.

\begin{remark}[Pressure in the enstrophy equation
vs.\ the stretching equation]
\label{rem:degree4}
The ratio $|\int\omega\cdot P\omega|/\mathcal{W}$
grows with Re (from $0.12$ at Re$\,=100$ to $2.58$
at Re$\,=1000$) because $\int\omega\cdot P\omega$
is degree~4 in $u$ while $\mathcal{W}$ is degree~3.
However, the \emph{evolution} of $\mathcal{W}$
(Proposition~\ref{prop:W-evol} below) reveals that
$\int\omega\cdot P\omega$ competes with $\|S\omega\|^2$
(also degree~4), not with $\mathcal{W}$ (degree~3).
The ratio $|\int\omega\cdot P\omega|/\|S\omega\|^2
\approx 0.32$ at Re$\,=1000$ --- pressure reduces the
self-amplification by a third, but does not exceed it.

The N14--N15 data remains significant for the
\emph{enstrophy equation}, where the pressure Hessian
enters indirectly through the time-evolution of
$\mathcal{W}$: the 32\% reduction of
$d\mathcal{W}/dt$ translates to a slower growth rate
for $\mathcal{W}$, which in turn limits the enstrophy
growth $d\mathcal{E}/dt = 2\mathcal{W}
- 2\nu\|\nabla\omega\|^2$.
\end{remark}

\begin{remark}[Threshold Reynolds number]
\label{rem:threshold}
The data suggests a critical
$\mathrm{Re}_c\approx 200$ above which
$\int\omega\cdot P\omega < 0$ at peak enstrophy
(pressure opposes stretching).  Below $\mathrm{Re}_c$,
viscosity alone controls the enstrophy.
Above $\mathrm{Re}_c$, pressure provides an
additional stabilization mechanism whose relative
contribution grows with Re.
\end{remark}


%======================================================================
\section{The Stabilization Conjecture}\label{sec:conjecture}
%======================================================================

\begin{proposition}[Evolution of vortex stretching]
\label{prop:W-evol}
For smooth solutions of~\eqref{eq:NS}:
\begin{equation}\label{eq:W-evol}
\frac{d\mathcal{W}}{dt} = \|S\omega\|^2
+ \int\omega\cdot P\omega\,dx
- 2\nu\int S_{ij}(\partial_k\omega_i)
(\partial_k\omega_j)\,dx.
\end{equation}
The three terms:
\begin{itemize}[leftmargin=2em]
\item $\|S\omega\|^2 > 0$: self-amplification of
  stretching (always positive, degree~4).
\item $\int\omega\cdot P\omega$: pressure Hessian
  contribution (negative for developed turbulence,
  degree~4).
\item $-2\nu\int S:(\nabla\omega)^T(\nabla\omega)$:
  viscous strain--vorticity-gradient interaction
  (negative when $\nabla\omega$ aligns with the
  compressive strain direction).
\end{itemize}
\end{proposition}

\begin{proof}
$\mathcal{W} = \int\omega\cdot S\omega\,dx$.
Transport terms: $(u\cdot\nabla)(\omega\cdot S\omega)
= 2((u\cdot\nabla)\omega)\cdot S\omega
+ \omega\cdot((u\cdot\nabla)S)\omega$,
integrating to zero on $\T^3$ by $\nabla\cdot u=0$.
Stretching from $\partial_t\omega$: $2\int(S\omega)
\cdot S\omega = 2\|S\omega\|^2$.
Stretching from $\partial_t S$:
$-\int\omega\cdot S^2\omega = -\|S\omega\|^2$.
Net: $+\|S\omega\|^2$.
Rotation: $\int\omega\cdot W^2\omega = 0$
(since $W\omega = 0$ for antisymmetric $W$
with associated vector $\omega/2$).
Pressure: $+\int\omega\cdot P\omega$ from the
strain equation (trace-free projection terms
cancel).

The raw viscous contribution is
$\nu(2\int\Delta\omega\cdot S\omega
+ \int\omega\cdot\Delta S\,\omega)$.
By Miller's identity~\eqref{eq:miller}:
$\int\omega\cdot\Delta S\,\omega
= -\langle\Delta S,\omega\otimes\omega\rangle = 0$.
Integration by parts on the remaining term:
$2\nu\int\Delta\omega\cdot S\omega
= -2\nu\int(\partial_k\omega_i)
\partial_k(S_{ij}\omega_j)$
$= -2\nu\int S_{ij}(\partial_k\omega_i)(\partial_k\omega_j)
- 2\nu\int(\partial_k\omega_i)(\partial_k S_{ij})\omega_j$.
The cross term vanishes by the same integration
by parts that yields Miller's identity.

Verified numerically to $0.78\%$ accuracy via
5-point finite difference stencil at
$\Delta t = 10^{-3}$, and Miller's identity
$\int\omega\cdot\Delta S\,\omega = 0$ verified
to machine precision.
\end{proof}

\begin{remark}[Role of Miller's identity]
\label{rem:miller-role}
Miller's identity~\eqref{eq:miller} plays a
simplifying role: it eliminates one of the two
viscous terms, leaving only the
strain--vorticity-gradient interaction.
The surviving term
$-2\nu\int S_{ij}(\partial_k\omega_i)
(\partial_k\omega_j)$ is a quadratic form in
$\nabla\omega$ weighted by the strain tensor.
Since $S$ is trace-free with both positive
(extensional) and negative (compressive)
eigenvalues, this form is \emph{not}
sign-definite.  However, in developed turbulence,
vorticity gradients are concentrated in the
compressive strain direction (perpendicular to
vortex tubes), making the quadratic form
negative --- i.e., viscosity opposes stretching.
\end{remark}

\begin{remark}[What the $\mathcal{W}$-evolution
reveals]
\label{rem:W-evol-meaning}
At Re$\,=1000$ (peak enstrophy):
$\|S\omega\|^2 = 1.54\times 10^4$ (self-amplification),
$\int\omega\cdot P\omega = -4.95\times 10^3$
(32\% reduction by pressure),
$\text{viscous} = -4.12\times 10^3$
(27\% reduction by viscosity).
Together, pressure and viscosity reduce the
self-amplification by 59\%, but do not reverse it.

The self-limitation of $\mathcal{W}(t)$ observed in
N11--N12 arises from the full feedback loop:
$\mathcal{W}$ drives enstrophy growth via
$d\mathcal{E}/dt = 2\mathcal{W}
-2\nu\|\nabla\omega\|^2$; enstrophy growth
increases dissipation (which scales as
$\|\nabla\omega\|^2$); eventually dissipation
catches $\mathcal{W}$, enstrophy peaks, and the
vorticity concentration that feeds $\|S\omega\|^2$
subsides.
\end{remark}

\subsection{The stabilization conjecture, revised}

\begin{conjecture}[Pressure stabilization]
\label{conj:pressure}
For smooth solutions of~\eqref{eq:NS} in the
developed turbulent regime
($\mathrm{Re}_{\rm eff}\gtrsim 200$):
\begin{equation}\label{eq:P-conj}
\int_{\T^3}\omega\cdot P\omega\,dx \leq 0,
\end{equation}
and the ratio $|\int\omega\cdot P\omega|
/\|S\omega\|^2$ is bounded below by a positive
constant independent of $\|\omega\|$:
\begin{equation}\label{eq:P-conj-ratio}
\frac{|\int\omega\cdot P\omega|}{\|S\omega\|^2}
\geq c_P > 0.
\end{equation}
Numerically, $c_P\approx 0.3$--$0.5$ across
Re$\,=200$--$1000$ (N14).
\end{conjecture}

\begin{remark}[What the conjecture gives and
what it does not]
\label{rem:conj-consequence}
If Conjecture~\ref{conj:pressure} holds with
$c_P\approx 0.3$, the $\mathcal{W}$-evolution
becomes
$d\mathcal{W}/dt \leq (1-c_P)\|S\omega\|^2
+ \text{viscous}$:
a 30\% reduction of the self-amplification rate.
This does \emph{not} directly give regularity
(since $(1-c_P)\|S\omega\|^2>0$ still drives
$\mathcal{W}$ upward), but it quantifies the
pressure contribution to the self-limitation
mechanism observed in N11--N13.

Regularity would require the full feedback loop
(enstrophy $\to$ dissipation $\to$ enstrophy
saturation) to be controlled, which involves
the interplay between the $\mathcal{W}$-evolution
and the enstrophy equation.  The pressure Hessian
is one of three stabilizing mechanisms
(alongside viscous dissipation and eigenvalue
selection) that together limit enstrophy growth.
Whether their combined effect precludes blow-up
remains the central open question.
\end{remark}


%======================================================================
\section{Connection to Regularity}\label{sec:regularity}
%======================================================================

\begin{proposition}[Depleted enstrophy equation]
\label{prop:depleted}
If $\rho(u(t))\leq\rho_{\max}<1$ for all $t\geq 0$,
then:
\begin{equation}\label{eq:depleted-enst}
\frac{1}{2}\frac{d}{dt}\|\omega\|^2
\leq \rho_{\max}\int|\omega|^2|S|\,dx
- \nu\|\nabla\omega\|^2.
\end{equation}
The effective nonlinearity is reduced by the factor
$\rho_{\max}$.  With $\rho_{\max}\approx 0.19$
(as measured in Section~\ref{sec:numerics}), the
stretching is $5\times$ weaker than the worst case.
\end{proposition}

\begin{proof}
$\mathcal{W} = \rho\int|\omega|^2|S|\,dx
\leq\rho_{\max}\int|\omega|^2|S|\,dx$.
\end{proof}

\begin{remark}[The degree-3 wall and the
three stabilizing mechanisms]
\label{rem:wall}
The bound $\rho_{\max}<1$ reduces the constant
but not the exponent: $\int|\omega|^2|S|\,dx$
is still degree~3.  Regularity does not follow
from $\rho<1$ alone.

The $\mathcal{W}$-evolution
(Proposition~\ref{prop:W-evol}) shows that
$d\mathcal{W}/dt = \|S\omega\|^2
+ \int\omega\cdot P\omega + \text{viscous}$,
where both $\|S\omega\|^2$ and
$\int\omega\cdot P\omega$ are degree~4.
At Re$\,=1000$, pressure reduces
self-amplification by $\sim 32\%$ and viscosity
by $\sim 27\%$, but their combined $59\%$ reduction
does not reverse the sign of $d\mathcal{W}/dt$.

What actually limits enstrophy growth is the
\emph{feedback loop}: stretching drives enstrophy,
enstrophy drives dissipation, and dissipation
eventually catches stretching.  The pressure
Hessian and eigenvalue selection slow this loop,
while viscosity provides the ultimate saturation.
Whether this feedback always prevents blow-up
remains the central open question.
\end{remark}

\begin{remark}[Comparison with existing regularity
criteria]
\label{rem:comparison}
The eigenvalue selection bound $\rho<1$ is
equivalent to requiring that vorticity is not
perfectly aligned with the most extensional
eigenvector --- a condition that holds at
\emph{every} point of \emph{every} smooth NS
solution, since perfect alignment at every point
would require $\omega$ to be an eigenvector of
$S$ everywhere, which is generically false.
The nontrivial content is the \emph{quantitative}
bound: $\rho\leq 0.19+\varepsilon$ rather than
merely $\rho<1$.

Miller~\cite{Miller24} obtains regularity criteria
involving the alignment of advection with the
strain self-interaction; the present framework
is complementary, working with the
$\hat\omega\cdot S\hat\omega$ projection directly.
\end{remark}


%======================================================================
\section{Summary}\label{sec:summary}
%======================================================================

\begin{center}
\renewcommand{\arraystretch}{1.2}
\begin{tabular}{@{}lll@{}}
\toprule
\textbf{Result} & \textbf{Status} & \textbf{Section} \\
\midrule
$\mathcal{W}=\int|\omega|^2\alpha\,dx$
  (stretching decomposition)
  & Proved & \ref{sec:decomp} \\
$\rho^*=1/\sqrt{26}$ (RE fixed point)
  & Proved & \ref{sec:RE} \\
RE stability of intermediate alignment
  & Proved & \ref{sec:RE} \\
$\alpha$-evolution equation
  & Proved & \ref{sec:GG} \\
$\int\omega\cdot P\omega = \int p\,\tr((\nabla\omega)^2)$
  & Proved & \ref{sec:pressure} \\
$\rho\approx 0.19$ (10 seeds, saturation by $t=2$)
  & Numerical & \ref{sec:numerics} \\
$P_\omega<0$ for developed turbulence
  & Numerical & \ref{sec:numerics} \\
$|P_\omega|$ comparable to RE self-damping
  & Numerical & \ref{sec:numerics} \\
$\int\omega\cdot P\omega = \int p\,\tr((\nabla\omega)^2)$
  verified to machine precision (N9)
  & Numerical & \ref{sec:numerics} \\
Sign transition: $\int\omega\cdot P\omega>0$ (early)
  $\to<0$ (developed) (N10)
  & Numerical & \ref{sec:numerics} \\
Near-cancellation: sub-terms $O(10^3)$,
  difference $O(1)$ (N10)
  & Numerical & \ref{sec:numerics} \\
$\mathcal{W}$ peaks before enstrophy;
  self-limiting growth rate (N11--N12)
  & Numerical & \ref{sec:numerics} \\
$\int\omega\cdot P\omega/\mathcal{W}\to -0.9$
  at peak enstrophy (N13)
  & Numerical & \ref{sec:numerics} \\
$|\!\int\omega\cdot P\omega|/\mathcal{W}$
  grows with Re: $0.1\to 2.6$ (N14)
  & Numerical & \ref{sec:numerics} \\
At Re$\,=1000$: ratio $<-2$ at all scales
  $k>2$ to $k>32$ (N15)
  & Numerical & \ref{sec:numerics} \\
$\rho\approx 0.14\sqrt{1-\mathcal{C}^{(1)}}$
  (surgical)
  & Numerical & \ref{sec:numerics} \\
$d\mathcal{W}/dt = \|S\omega\|^2
  +\int\omega\cdot P\omega+\text{visc}$
  & Proved & \ref{sec:conjecture} \\
Pressure$+$viscosity reduce $\|S\omega\|^2$
  by $59\%$ at Re$\,=1000$
  & Numerical & \ref{sec:conjecture} \\
$\int\omega\cdot P\omega\leq 0$,
  $|\cdot|/\|S\omega\|^2\geq c_P$
  & Conjecture & \ref{sec:conjecture} \\
\bottomrule
\end{tabular}
\end{center}


%======================================================================
\section*{Acknowledgments}
%======================================================================

AI tools were used as interactive assistants for
drafting and mathematical exploration.
Simulations were performed on an NVIDIA RTX~5090.

%======================================================================
\begin{thebibliography}{99}

\bibitem{AKKG87}
W.\,T.~Ashurst, A.\,R.~Kerstein, R.\,M.~Kerr,
and C.\,H.~Gibson,
Alignment of vorticity and scalar gradient with strain
rate in simulated Navier--Stokes turbulence,
\emph{Phys.\ Fluids}~\textbf{30} (1987), 2343--2353.

\bibitem{Vieillefosse82}
P.~Vieillefosse,
Local interaction between vorticity and shear in a
perfect incompressible fluid,
\emph{J.\ Physique}~\textbf{43} (1982), 837--842.

\bibitem{Cantwell92}
B.\,J.~Cantwell,
Exact solution of a restricted Euler equation for the
velocity gradient tensor,
\emph{Phys.\ Fluids A}~\textbf{4} (1992), 782--793.

\bibitem{GalantiGibbon97}
B.~Galanti, J.\,D.~Gibbon, and M.~Heritage,
Vorticity alignment results for the three-dimensional
Euler and Navier--Stokes equations,
\emph{Nonlinearity}~\textbf{10} (1997), 1675--1694.

\bibitem{Miller24}
E.~Miller,
On the interaction of strain and vorticity for
solutions of the Navier--Stokes equation,
\emph{arXiv preprint} arXiv:2407.02691, 2024.

\bibitem{Tsinober92}
A.~Tsinober, E.~Kit, and T.~Dracos,
Experimental investigation of the field of velocity
gradients in turbulent flows,
\emph{J.\ Fluid Mech.}~\textbf{242} (1992), 169--192.

\bibitem{JimenezWray98}
J.~Jim\'enez and A.\,A.~Wray,
On the characteristics of vortex filaments in
isotropic turbulence,
\emph{J.\ Fluid Mech.}~\textbf{373} (1998), 255--285.

\bibitem{Fontes2026spread}
R.~Fontes,
Angular spreading of enstrophy in the 3D Navier--Stokes
equations,
\emph{Preprint}, 2026.

\bibitem{Fontes2026burgers}
R.~Fontes,
Pressure stabilization near the Burgers vortex:
the neutral point of eigenvalue selection,
\emph{Preprint}, 2026.

\end{thebibliography}

\end{document}
